Kiến trúc tổng thể của hệ thống được tổ chức theo mô hình phân tầng gồm tầng cảm biến,
tầng thiết bị nhúng ESP32, tầng truyền thông MQTT, tầng backend/API, tầng cơ sở dữ liệu
và tầng dashboard web. Cách tổ chức này giúp phân tách rõ trách nhiệm giữa khâu thu
thập dữ liệu, khâu truyền nhận qua mạng, khâu lưu trữ lịch sử và khâu giám sát, điều
khiển ở phía người dùng.

\section{Luồng dữ liệu}

Về mặt chức năng, hệ thống gồm hai chiều luồng chính: luồng giám sát từ thiết bị lên máy
chủ và luồng điều khiển từ người dùng xuống thiết bị. Luồng giám sát được mô tả như sau:

\begin{center}
\texttt{Sensor -> ESP32/Firmware -> MQTT -> Backend -> Database -> API -> Web Dashboard}
\end{center}

Trong luồng này, cảm biến cung cấp dữ liệu thô cho firmware trên ESP32; firmware thực
hiện đọc mẫu, kiểm tra hợp lệ, tổng hợp trạng thái IAQ và publish telemetry qua MQTT.
Backend tiếp nhận payload, kiểm tra theo contract chung, cập nhật trạng thái hiện thời,
lưu lịch sử vào MySQL và cung cấp API cho dashboard truy xuất.

Ở chiều ngược lại, luồng điều khiển được tổ chức theo hướng:

\begin{center}
\texttt{Web Dashboard -> API -> Backend -> MQTT -> ESP32 -> Device/LED}
\end{center}

Theo đó, thao tác của người dùng trên dashboard không tác động trực tiếp đến phần cứng,
mà đi qua backend và topic điều khiển MQTT trước khi tới firmware. Cách tổ chức này cho
phép chuẩn hóa lệnh điều khiển, lưu vết thao tác và áp dụng các ràng buộc an toàn trước
khi thiết bị thực thi.

Trong nguyên mẫu hiện tại, hai chiều luồng trên đã được hiện thực ở mức cơ bản bởi
firmware RTOS, backend, cơ sở dữ liệu và dashboard web.

\section{Luận điểm thiết kế hệ thống}

Luận điểm thiết kế cốt lõi của đề tài là ưu tiên xử lý cục bộ tại thiết bị trước, sau
đó mới mở rộng sang giám sát và điều khiển từ xa. Điều này có nghĩa là ESP32 vẫn phải
duy trì được các chức năng cơ bản như đọc cảm biến, đánh giá IAQ và kích hoạt cảnh báo
cục bộ ngay cả khi Wi-Fi, broker MQTT hoặc server tạm thời không khả dụng. Tầng IoT,
backend và dashboard vì vậy chỉ đóng vai trò mở rộng khả năng quan sát, lưu trữ lịch sử
và tương tác từ xa.

Từ luận điểm này, các quyết định thiết kế được sắp xếp theo thứ tự ưu tiên sau: dữ
liệu cảm biến phải được kiểm tra trước khi dùng cho điều khiển; cảnh báo cục bộ phải có
khả năng vận hành độc lập với server; các lệnh từ dashboard phải đi qua backend và
contract truyền thông thống nhất; còn cơ sở dữ liệu phục vụ chủ yếu cho lưu trữ và phân
tích lịch sử thay vì là điều kiện bắt buộc để thiết bị phản ứng tức thời.

\section{Phân rã module}

Trên cơ sở kiến trúc phân tầng, hệ thống được phân rã thành các module với trách nhiệm
xử lý tương đối độc lập. Việc phân rã này nhằm làm rõ ranh giới chức năng giữa các khối
thành phần và quan hệ trao đổi dữ liệu giữa chúng.

\begin{table}[H]
    \centering
    \caption{Phân rã module theo trách nhiệm kiến trúc}
    \begin{tabularx}{\textwidth}{>{\bfseries}l X X}
        \toprule
        Module & Chức năng chính & Quan hệ với các module khác \\
        \midrule
        \texttt{rtos} & Đọc cảm biến, xử lý dữ liệu cục bộ, đánh giá IAQ, điều khiển thiết bị và giao tiếp MQTT ở phía ESP32 & Nhận cấu hình và lệnh điều khiển từ backend qua MQTT; gửi telemetry và trạng thái thiết bị lên các tầng phía trên \\
        \texttt{common} & Định nghĩa topic MQTT, cấu trúc payload JSON và contract dữ liệu dùng chung & Là lớp giao ước giữa firmware, backend, cơ sở dữ liệu và frontend \\
        \texttt{iot\_backend/backend} & Tiếp nhận dữ liệu MQTT, kiểm tra payload, duy trì trạng thái hiện thời, cung cấp API và phát lệnh điều khiển & Kết nối giữa \texttt{rtos}, \texttt{database} và \texttt{frontend}; là điểm trung gian cho cả giám sát lẫn điều khiển \\
        \texttt{iot\_backend/database} & Lưu trữ telemetry lịch sử và nhật ký lệnh điều khiển & Nhận dữ liệu ghi từ backend và cung cấp nguồn dữ liệu lịch sử cho API/dashboard \\
        \texttt{iot\_backend/frontend} & Hiển thị dữ liệu realtime, dữ liệu lịch sử, trạng thái thiết bị và giao diện điều khiển người dùng & Khai thác dữ liệu từ backend qua API; không giao tiếp trực tiếp với phần cứng \\
        \texttt{integration\_test} & Hỗ trợ kiểm thử tích hợp luồng dữ liệu và luồng điều khiển ở mức hệ thống & Kiểm chứng sự tương thích giữa contract, backend và các luồng MQTT trước khi kiểm thử end-to-end thực tế \\
        \bottomrule
    \end{tabularx}
\end{table}

\section{Nguyên tắc thiết kế}

Ở mức triển khai, các tầng và module được tách theo trách nhiệm xử lý. Nhóm driver chỉ
đảm nhiệm giao tiếp với cảm biến và kiểm tra trạng thái giao thức. Lớp gom mẫu dữ liệu
chuẩn hóa thông tin đo về một cấu trúc chung để phục vụ các bước xử lý tiếp theo. Khối
đánh giá IAQ chỉ tập trung vào logic suy diễn trạng thái môi trường và ngưỡng điều
khiển, không phụ thuộc trực tiếp vào phần cứng đầu ra. Khối điều khiển đầu ra chịu trách
nhiệm áp dụng trạng thái thiết bị theo kết quả đánh giá hoặc theo lệnh điều khiển hợp
lệ.

Trong phiên bản RTOS, các thành phần này được tách thành nhiều tác vụ và trao đổi dữ
liệu thông qua queue, mutex và semaphore. Ở phía server, backend và dashboard không làm
việc trực tiếp với phần cứng mà giao tiếp với thiết bị thông qua MQTT và API dựa trên
cùng một contract dữ liệu. Cách tổ chức này giúp hệ thống nhất quán giữa tầng nhúng và
tầng ứng dụng, đồng thời thuận lợi cho việc mở rộng về sau.
